\subsection{仿射 Springer 纤维的更多例子}\label{ss:aff examples}
本小节再给出几个例子，展示仿射 Springer 纤维的丰富几何。导至这些几何描述的计算一概从略。

下面的所有例子中，仿射 Springer 纤维都是{\em 齐性}的：$\sX_{\gamma}$ 上多了一个 $\Gm$-作用，它来自 $\Gr_{G}$ 上的{\em 圈旋转}，即在素元 $t$ 上作伸缩。关于齐性仿射 Springer 纤维及其在表示论中的应用，参见 \cite{OY}。

\subsubsection{}\label{SL2 ram} 设 $G=\SL_{2}$ 与 $\gamma=\mat{0}{t^{m+1}}{t^{m}}{0}$。当 $m$ 为偶数时，$\sX_{\gamma}$ 恰由那些满足
\begin{equation*}
t^{m/2}\calO_{F}\oplus t^{m/2}\calO_{F}\subset\L\subset t^{-m/2}\calO_{F}\oplus t^{-m/2}\calO_{F}
\end{equation*}
的格 $\L$ 组成。此时 $\sX_{\gamma}$ 与 $\Gr_{G}$ 中对应于余权 $m\alpha^{\vee}/2$ 的 $L^{+}G$-轨道的闭包重合。

当 $m$ 为奇数时，$\sX_{\gamma}$ 恰由那些满足
\begin{equation*}
t^{(m+1)/2}\calO_{F}\oplus t^{(m-1)/2}\calO_{F}\subset\L\subset t^{-(m-1)/2}\calO_{F}\oplus t^{-(m+1)/2}\calO_{F}
\end{equation*}
的格 $\L$ 组成。此时改为考虑 $G'=\PGL_{2}$ 的仿射 Grassmann 簇 $\Gr_{G'}$，它把 $\Gr_{G}$ 作为其一个分支包含在内。$\sX_{\gamma}$ 可以等同于 $\Gr_{G'}$ 中对应于余权 $m\alpha^{\vee}/2$ 的 $L^{+}G'$-轨道的闭包。

无论哪种情形，都有 $\dim\sX_{\gamma}=m$。

\subsubsection{Lusztig--Smelt 例}\label{sss:LS} 设 $G=\GL_{n}$，$\gamma\in\frg(F)$ 的特征多项式为 $P(x)=x^{n}-t^{m}=0$，其中 $(m,n)=1$。若格 $\L\subset F^{n}$ 在 $\gamma$ 下稳定，则它带有环 $R=\calO_{F}[X]/(X^{n}-t^{m})$ 的作用。令 $K=\Frac(R)$，则 $K= k\lr{s}$，其中 $x=s^{m}$、$t=s^{n}$。于是 $R$ 在 $K$ 中的整闭包是 $\tilR:=k\tl{s}$。$\gamma$ 在 $F^{n}$ 上的作用使之成为一维 $K$-向量空间。固定一个 $K$-线性同构 $F^{n}\cong K$，在此同构下，$F^{n}$ 中在 $R$ 下稳定的格就只是{\em 分式 $R$-理想}，即 $K=\Frac(R)$ 的有限生成 $R$-子模。于是 $\sX_{\gamma}$ 可以等同于分式 $R$-理想的模空间。

中心化子 $G_{\gamma}(F)$ 就是 $K^{\times}$，它通过乘法作用在分式 $R$-理想集合上。此作用显然经由 $K^{\times}/R^{\times}$ 分解，而后者正是局部 Picard 群概形 $P_{\gamma}$ 的 $k$ 点群。

存在 $\Gm$ 在 $K$ 上的作用（由域自同构给出），它在 $s$ 上作伸缩（于是 $s^{i}$ 在此作用下有权 $i$）。此作用诱导 $\Gm$ 在 $\sX_{\gamma}$ 上的作用。$\Gm$ 在 $\sX_{\gamma}$ 上的不动点对应由 $s$ 的单项式生成的分式理想。更精确地说，若分式 $R$-理想 $\L\subset K=k\lr{s}$ 被 $\Gm$ 固定，则定义 $M_{\L}=\{i\in\ZZ|s^{i}\in \L\}$；它是 $\ZZ$ 的子集，且在加 $m$ 与加 $n$ 下稳定，因为 $\L$ 是 $R=k\tl{s^{m},s^{n}}$-模。于是 $M_{\L}\subset\ZZ$ 是幺半群 $A_{m,n}:=\ZZ_{\geq 0}m+\ZZ_{\geq0}n\subset\ZZ_{\geq 0}$ 上的有限生成模。对应 $\L\mapsto M_{\L}$ 给出双射
\begin{equation*}
\sX_{\gamma}^{\Gm}(k)\isom\{\mbox{$A_{m,n}$-子模 $M\subset \ZZ$}\}.
\end{equation*}
$\ZZ$ 的任一 $A_{m,n}$-子模都包含所有充分大的整数。因此任意两个这样的 $A_{m,n}$-模 $M$ 与 $M'$ 只差有限多个元素，从而可以定义 $[M:M']=\#(M\bs M')-\#(M'\bs M)$。固定任意 $i\in\ZZ$，满足 $[M:\ZZ_{\geq0}]=i$ 的不动点共有 $\frac{1}{n+m}\binom{n+m}{n}$ 个。对与 $A_{m,n}$-模 $M$ 对应的不动点 $p_{M}$，考虑子簇 $C_{M}=\{p\in\sX_{\gamma}|\lim_{\Gm\ni z\to0}z\cdot p=p_{M}\}$。则 $C_{M}$ 同构于仿射空间，其维数可按 $M$ 组合地表达。诸胞腔 $C_{M}$ 给出 $\sX_{\gamma}$ 的分层。由此可以计算 $\sX_{\gamma}$ 各连通分支的 Poincar\'e 多项式。更多细节以及 $\sY_{\gamma}$ 的类似图景，参见 \cite{LS}。

\subsubsection{} 我们在 \S\ref{sss:LS} 的一个特例中考察 $\sY_{\gamma}$ 的几何。取 $G=\GL_{3}$ 与 $\gamma^{3}-t^{2}=0$。如前引入变量 $s$，使 $t=s^{3}$、$\gamma=s^{2}$，则 $R=k\tl{s^{2},s^{3}}$，其分式域为 $K=k\lr{s}$。仿射 Springer 纤维 $\sY_{\gamma}$ 分类分式 $R$-理想的链 $\L_{0}\subset \L_{1}\subset \L_{2}\subset s^{-3}\L_{0}$。考虑分支 $\sY^{0}_{\gamma}\subset\sY_{\gamma}$，它分类如上并满足 $[\L_{i}:k\tl{s}]=i$（$0\leq i\leq 2$）的链 $\{\L_{i}\}$。

先研究 $\sY^{0}_{\gamma}$ 上的 $\Gm$-不动点。对每个被 $\Gm$ 固定的 $R$-子模 $k\lr{s}$，用一列整数记其相应的幺半群 $A_{3,2}$-模（见 \S\ref{sss:LS}）。例如 $(0,1,2,\cdots)$ 表示标准格 $k\tl{s}$。注意任一被 $\Gm$ 固定的 $R$-分式理想对应的数列，要么是连续的 $(n,n+1,n+2,\cdots)$，要么在第二处至多有一个缺口 $(n,n+2,n+3,\cdots)$。$\sY^{0}_{\gamma}$ 中有如下四个不动点：
\begin{enumerate}
\item $q: (0,1,2,\cdots)\subset (-1,0,1,\cdots)\subset (-2,-1,0,\cdots)$；
\item $p_{0}: (-1,1,2,\cdots)\subset (-1,0,1,\cdots)\subset (-2,-1,0,\cdots)$；
\item $p_{1}: (0,1,2,\cdots)\subset (-2,0,1,\cdots)\subset (-2,-1,0,\cdots)$；
\item $p_{2}:(0,1,2,\cdots)\subset (-1,0,1,\cdots)\subset (-3,-1,0,\cdots)$。
\end{enumerate}
我们的编号方式是：$p_{i}$ 由 $q$ 改变第 $i$ 个格 $\L_{i}$ 得到。

于是 $\sY^{0}_{\gamma}$ 是三个不可约分支之并 $C_{0}\cup C_{1}\cup C_{2}$，每个分支 $C_{i}$ 同构于 $\PP^{1}$。三者都过 $q$，且 $q$ 是任意两个分支唯一的交点。对 $i=0,1,2$ 有 $p_{i}\in C_{i}$。

如 \S\ref{sss:GLn parahoric} 所见，$G$ 的仿射部分旗簇可以自然地用子集 $J\subset \{0,1,2\}$ 编号。令 $\sX_{J,\gamma}$ 为 $\gamma$ 在 $\Fl_{J}$ 中的仿射 Springer 纤维。在投影 $\sY_{\gamma}\to \sX_{J,\gamma}$ 下，$i\notin J$ 的曲线 $C_{i}$ 缩成一点，其余的曲线同构地映到其像上。

\subsubsection{}\label{ex:Sp} 设 $G=\Sp(V)$，其中 $V$ 是 $k$ 上维数为 $2n$ 的辛空间，并假定 $\chark\neq2$。固定分解 $V=U\oplus U^{*}$ 为 $V$ 的 Lagrange 子空间，使辛形式限制在 $U\times U^{*}$ 上为自然配对。考虑 $\gamma=\mat{0}{tX}{Y}{0}\in\frg(F)$，其中 $X\in\Sym^{2}(U)$（视为自伴随映射 $\xi:U^{*}\to U$），$Y\in\Sym^{2}(U^{*})$（视为自伴随映射 $\eta:U\to U^{*}$）。$\gamma$ 正则半单的条件等价于：(1) $\xi$ 与 $\eta$ 都是同构；(2) $\xi\eta\in\GL(U)$ 正则半单，或等价地 $\eta\xi\in\GL(U^{*})$ 正则半单。

仿射 Springer 纤维 $\sX_{\gamma}$ 分类在 $\gamma$ 下稳定的自对偶格 $\L\subset V\otimes_{k}F$（见 \S\ref{sss:ASF Gr GLn}）。考虑 $\Gm$ 在 $V\otimes_{k}F=U\otimes_{k}F\oplus U^{*}\otimes_{k}F$ 上的作用，使 $Ut^{i}$ 的权为 $2i-1$、$U^{*}t^{i}$ 的权为 $2i$。此作用诱导 $\Gm$ 在 $\Gr_{G}$ 与 $\sX_{\gamma}$ 上的作用。先考虑不动点 $\Gr_{G}^{\Gm}$。格 $\L\in\Gr_{G}$ 被 $\Gm$ 固定，当且仅当它是
\begin{equation*}
\Span\{(A_{i}\oplus B_{i})t^{i}; i\in\ZZ \}
\end{equation*}
的 $t$-进完备化，其中 $\cdots \subset A_{-1}\subset A_{0}\subset A_{1}\subset \cdots \subset U$ 是 $U$ 的滤过，满足 $N\ll 0$ 时 $A_{N}=0$、$N\gg 0$ 时 $A_{N}=U$；而 $\cdots \subset B_{-1}\subset B_{0}\subset B_{1}\subset \cdots \subset U^{*}$ 是 $U^{*}$ 的类似滤过，并且在对偶配对下成立
\begin{equation}\label{AB dual}
B_{i-1}=A^{\bot}_{-i}, \forall i\in\ZZ.
\end{equation}
最后一个条件反映了 $\L$ 在辛形式下自对偶这一事实。

格 $\L\in\sX_{\gamma}^{\Gm}$ 于是由 $U$ 的滤过 $A_{\bu}$ 与 $U^{*}$ 的滤过 $B_{\bu}$ 决定，两者在 \eqref{AB dual} 的意义下对偶，并满足附加条件
\begin{equation*}
\eta A_{i}\subset B_{i}; \quad \xi B_{i}\subset A_{i+1}, \quad\forall i\in\ZZ.
\end{equation*}
把上述数据概括为下图
\begin{equation*}
\xymatrix{ A_{-1}\ar@{^{(}->}[rr]\ar[dr]^{\eta} && A_{0}\ar@{^{(}->}[rr]\ar[dr]^{\eta} && A_{1}\ar[dr]^{\eta}& \cdots\\
\cdots & A^{\bot}_{0}\ar@{^{(}->}[rr]\ar[ur]^{\xi} && A^{\bot}_{-1}\ar@{^{(}->}[rr]\ar[ur]^{\xi} && A_{-2}^{\bot}}
\end{equation*}

例如当 $\dim V=4$（从而 $\dim U=2$）时有两种可能。第一种是 $A_{-1}=0\subset A_{0}=U$，它对应标准格 $V\otimes_{k}\calO_{F}$。第二种是 $A_{-1}=0\subset A_{0}\subset A_{1}=U$，其中 $A_{0}\subset U$ 是一条满足 $\xi A^{\bot}_{0}\subset A_{0}$ 的直线，即 $A^{\bot}_{0}\subset U^{*}$ 是二次形式 $X$ 下的迷向直线。这样的直线 $A^{\bot}_{0}$ 有两条，给出另外两个 $\Gm$-不动点。因此 $\sX_{\gamma}^{\Gm}$ 由 3 个点组成。此时 $\sX_{\gamma}=C_{1}\cup C_{2}$，其中 $C_{i}\cong\PP^{1}$，两个分支交于标准格，且每个 $C_{i}$ 含有其余 $\Gm$-不动点之一。

\subsubsection{Bernstein--Kazhdan 例}\label{sss:BK example} 在 \cite[附录]{KL} 中，Bernstein 与 Kazhdan 给出了仿射 Springer 纤维的不可约分支不是有理簇的第一个例子。沿用 \S\ref{ex:Sp} 的记号。设 $\Fl_{\bP}$ 是 $G=\Sp(V)$ 的仿射部分旗簇，它分类满足 $\L^{\vee}=\L$ 与 $\L'^{\vee}=t^{-1}\L$ 的格对 $\L'\subset\L$。

取 \S\ref{ex:Sp} 中的 $\gamma$。同一个 $\Gm$ 作用在 $\sX_{\bP,\gamma}$ 上，而不动点 $\L'\subset \L$ 可以用两对滤过 $(A_{\bu}, B_{\bu})$ 与 $(A'_{\bu}, B'_{\bu})$ 来描述：其中 $(A_{\bu}, B_{\bu})$ 是 \S\ref{ex:Sp} 所述那种 $U$ 与 $U^{*}$ 的滤过，$(A'_{\bu}, B'_{\bu})$ 类似，只是把 \eqref{AB dual} 换成 $B'_{i}=A_{-i}'^{\bot}$。此外，包含 $\L'\subset \L$ 等价于对一切 $i$ 有 $A'_{i}\subset A_{i}$ 与 $B'_{i}\subset B_{i}$。

例如取 $\dim V=6$，并固定滤过的维数：
\begin{equation*}
A_{-1}=0, \quad \dim A_{0}=2, \quad A_{1}=U;\quad  \dim A'_{0}=1, \quad A'_{1}=U.
\end{equation*}
这样的滤过组 $(A_{\bu}, B_{\bu}; A'_{\bu}, B'_{\bu})$ 由满足 $\xi A^{\bot}_{0}\subset A_{0}$ 与 $\eta A'_{0}\subset A'^{\bot}_{0}$ 的完备旗 $0\subset A'_{0}\subset A_{0}\subset U$ 决定。换言之，$A^{\bot}_{0}$ 是 $U^{*}$ 中二次形式 $X$ 下的迷向直线，$A'_{0}$ 是 $U$ 中二次形式 $Y$ 下的迷向直线。数对 $(A^{\bot}_{0}, A'_{0})$ 决定 $Q(X)\times Q(Y)\subset \PP(U)\times\PP(U^{*})$ 中一点，这里 $Q(X)$ 与 $Q(Y)$ 分别是由 $X$ 与 $Y$ 定义的二次曲线。关联关系 $A'_{0}\subset A_{0}$ 在 $Q(X)\times Q(Y)\cong\PP^{1}\times \PP^{1}$ 中定义了一条双次数为 $(2,2)$ 的曲线，从而是一条亏格为 $1$ 的曲线。因此 $\sX^{0}_{\bP,\gamma}$ 的连通分支是一条亏格 $1$ 的曲线。考虑 $\sX_{\bP,\gamma}$ 中收缩到该曲线的点，并取其闭包 $Z$。可以证明 $\dim Z=3=\dim\sX_{\bP,\gamma}$。于是 $\sX_{\bP,\gamma}$ 含有一个非有理的不可约分支 $Z$。更多细节参见 \cite{KL} 的附录。

\subsubsection{「次正则」仿射 Springer 纤维}
当 $\sX_{\gamma}$ 或 $\sY_{\gamma}$ 是一维时，不妨称之为次正则仿射 Springer 纤维，以类比 \S\ref{sss:subreg} 中讨论的次正则 Springer 纤维。若 $\dim\sY_{\gamma}=1$，则它是 $\PP^{1}$ 之并，于是可以定义它的对偶图。在 \cite[命题 7.7]{KL} 中，$\Fl$ 中次正则仿射 Springer 纤维的对偶图得到了分类：它们几乎总是单连型群的扩张 Dynkin 图，只是在 $A$ 型时还可以是无穷长链（见例 \ref{P1 chain}）。
